% Irodov, user-supplied Arihant PDF, original pp.59--60 (PDF pp.68--69).
% Section 1.8: special relativity. Problems 361--370, full derivations.
% Numerical constants: c=2.99792458e8 m/s, g=9.81 m/s^2.

\exproblem{361}{59}
两相对论性粒子在实验室系中沿相互垂直的方向运动，
速率分别为 $v_1,v_2$。
求两粒子之间的相对速率。
\exsolution
取粒子1的实验室速度沿 $+x$ 方向，即
$\bm v_1=v_1\bm e_x$；粒子2沿 $+y$，
$\bm v_2=v_2\bm e_y$。
转到粒子1的瞬时静止惯性系，其相对实验室系
沿 $+x$ 以 $v_1$ 运动。速度分量变换为
\[
v'_{2x}=\frac{0-v_1}{1-0}=-v_1,\qquad
v'_{2y}=\frac{v_2\sqrt{1-v_1^2/c^2}}{1-0}
=v_2\sqrt{1-v_1^2/c^2}.
\]
两分量平方求和，得两粒子的相对速率
\[
\boxed{v_{\rm rel}=
\sqrt{v_1^2+v_2^2-\frac{v_1^2v_2^2}{c^2}}}.
\]
\excheck 当任一速率为零时退化为另一速率，
而当二者均小于 $c$ 时，结果仍小于 $c$，
因为
$1-v_{\rm rel}^2/c^2=
(1-v_1^2/c^2)(1-v_2^2/c^2)>0$。

\exproblem{362}{59}
一个不稳定粒子在参考系 $K'$ 中沿 $y'$ 轴以
速率 $v'$ 运动。$K'$ 相对 $K$ 沿 $+x$ 轴
以速率 $V$ 匀速运动，$x,x'$ 轴同向、
$y,y'$ 轴平行。粒子的固有寿命为
$\Delta t_0$，求在 $K$ 系中粒子实际走过的距离。
\exsolution
令 $\beta=V/c$、
$\gamma_V=1/\sqrt{1-\beta^2}$、
$\gamma_{v'}=1/\sqrt{1-v'^2/c^2}$。
在 $K'$ 中粒子的速度为
$(v_x',v_y')=(0,v')$。
反向速度变换给
\[
v_x=V,\qquad v_y=\frac{v'}{\gamma_V},
\quad
v_K=\sqrt{V^2+(1-\beta^2)v'^2}.
\]
粒子在 $K'$ 中的寿命为
$\Delta t'=\gamma_{v'}\Delta t_0$。
世界线在 $K'$ 中满足 $\Delta x'=0$，
因此洛伦兹时间变换得到
$\Delta t=\gamma_V\Delta t'=
\gamma_V\gamma_{v'}\Delta t_0$。
于是在 $K$ 中的路程
\[
\boxed{s=v_K\Delta t=
\Delta t_0\sqrt{
\frac{V^2+(1-\beta^2)v'^2}
{(1-\beta^2)(1-v'^2/c^2)}}}.
\]
\excheck 当 $V=0$ 时
$s=v'\gamma_{v'}\Delta t_0$；
当 $v'=0$ 时
$s=V\gamma_V\Delta t_0$，
均符合单一参考系的时间膨胀关系。

\exproblem{363}{59}
一粒子在 $K$ 系内运动，速率为 $v$，
速度矢量与 $+x$ 轴夹角为 $\theta$。
$K'$ 系相对 $K$ 沿 $+x$ 方向
匀速运动，速率为 $V$，
两系 $x,x'$ 轴同向。
求在 $K'$ 系测得的运动方向角 $\theta'$。
\exsolution
在 $K$ 中
$v_x=v\cos\theta$、
$v_y=v\sin\theta$。
令 $\beta=V/c$，
洛伦兹速度变换为
\[
v_x'=\frac{v\cos\theta-V}
{1-vV\cos\theta/c^2},\qquad
v_y'=\frac{v\sin\theta\sqrt{1-\beta^2}}
{1-vV\cos\theta/c^2}.
\]
二者相除，得到
\[
\boxed{\tan\theta'=
\frac{\sqrt{1-\beta^2}\sin\theta}
{\cos\theta-V/v}}.
\]
\excheck 该式直接给出正切值；
若分母改变符号，必须依据
$(v_x',v_y')$ 所在象限选取
$\theta'$，不可只用 $\arctan$ 的主值。

\exproblem{364}{59}
如原书图1.94，一根细杆 $AB$ 在 $K'$ 系中
平行于 $x'$ 轴，且沿 $y'$ 负方向
以速度 $v'$ 平动。
与此同时，$K'$ 系相对实验室 $K$
沿 $+x$ 方向匀速运动，速率为 $V$。
求在 $K$ 系中\emph{同一时刻}看到的杆与
$x$ 轴之间的倾角大小 $\theta$。
\exsolution
杆在 $K'$ 的一个共同时间上两端满足
$\Delta x'=L'$、$\Delta y'=0$；
因杆还沿 $y'$ 方向平动，
两个任意取样事件满足
$\Delta y'=-v'\Delta t'$。
在实验室 $K$ 的同时性条件 $\Delta t=0$ 下，
洛伦兹变换
$\Delta t=\gamma_V(\Delta t'+V\Delta x'/c^2)$
给出
\[
\Delta t'=-\frac{V\Delta x'}{c^2},
\quad
\Delta y'=\frac{Vv'}{c^2}\Delta x'.
\]
实验室的两端横向距离
$\Delta y=\Delta y'$，
而
$\Delta x=\gamma_V(\Delta x'+V\Delta t')
=\Delta x'/\gamma_V$。
故
\[
\boxed{\tan\theta=
\gamma_V\frac{Vv'}{c^2}
=\frac{Vv'}{c^2\sqrt{1-V^2/c^2}}}.
\]
\excheck 这里的倾斜来自\emph{同时性相对性}：
实验室同时测量的两端在 $K'$ 中并非同时，
而杆在两次 $K'$ 时刻之间有横向位移。

\exproblem{365}{59}
惯性系 $K'$ 相对 $K$ 以恒速 $V$
运动。粒子在 $K$ 中沿一条直线运动，
瞬时速率为 $v$、加速度模为 $w$。
求它在 $K'$ 系的瞬时加速度 $w'$：
(a) 粒子速度平行于 $\bm V$；
(b) 粒子速度垂直于 $\bm V$。
\exsolution
取 $K'$ 沿 $+x$ 方向移动，
$\beta=V/c$，并从速度变换求导。
(a) 粒子沿 $x$ 运动，有
$v'_x=(v-V)/(1-vV/c^2)$。
对实验室时间求导：
\[
\frac{\mathrm dv'_x}{\mathrm dt}
=\frac{(1-\beta^2)w}{(1-vV/c^2)^2},
\quad
\frac{\mathrm dt'}{\mathrm dt}
=\gamma_V(1-vV/c^2).
\]
二者相除：
\[
\boxed{w'_\parallel=
\frac{w(1-\beta^2)^{3/2}}
{(1-vV/c^2)^3}}.
\]
若 $\bm v$ 与 $\bm V$ 反向，则
沿 $+x$ 的有符号 $v$ 应取负值。
(b) 若粒子沿 $y$ 运动并在该方向加速，
则实验室 $v_x=0$ 且始终 $a_x=0$。
$K'$ 中横向速度 $v'_y=v_y/\gamma_V$，
同时 $\mathrm dt'/\mathrm dt=\gamma_V$，
所以
\[
\boxed{w'_\perp=w(1-\beta^2)}.
\]
\excheck (a) 的因子取决于速度方向相对 $V$ 的
有符号关系；(b) 是粒子运动和加速度都垂直于
两系相对运动方向的情形。

\exproblem{366}{59}
一枚理想火箭从地球静止起飞，
其在每个瞬时共动惯性系中测得的
固有加速度均为 $w'=10g$，
推进阶段持续地球参考系时间
$\tau=1.0$ 年。
求推进末速比光速少百分之多少，以及
地球系中飞行的距离。
\exsolution
设固有加速度 $a=w'=10g$，
在地球系中火箭沿一直线加速。
一维加速度的变换关系为
$a=\gamma^3\,\mathrm dv/\mathrm dt$，于是
\[
\frac{\mathrm dv}{\mathrm dt}
=a\left(1-\frac{v^2}{c^2}\right)^{3/2}.
\]
由 $v(0)=0$ 积分：
\[
\frac{v}{\sqrt{1-v^2/c^2}}=at.
\]
解为
\[
\boxed{v(t)=\frac{at}{\sqrt{1+(at/c)^2}}}.
\]
因此末速相对光速的百分差为
\[
\boxed{100\frac{c-v(\tau)}c
=100\left[1-\frac{a\tau/c}{\sqrt{1+(a\tau/c)^2}}\right]
\simeq0.47\%}.
\]
对 $v(t)$ 再积分得
\[
\boxed{x(\tau)=\frac{c^2}{a}
\left[\sqrt{1+(a\tau/c)^2}-1\right]
\simeq0.91\ \text{光年}}.
\]
\excheck 一年在本题是\emph{地球系时间}，
不能拿它直接代入火箭钟的固有时长关系。
这里忽略引力及实际推进剂消耗。

\exproblem{367}{59}
沿用上一题的火箭与恒定固有加速度
$w'=10g$、地球时间 $\tau=1.0$ 年。
求推进阶段火箭自身时钟记录的固有时间
$\tau_0$。
\exsolution
由上一题速度规律
$v(t)=at/\sqrt{1+(at/c)^2}$，
得到
\[
\sqrt{1-v(t)^2/c^2}
=\frac1{\sqrt{1+(at/c)^2}}.
\]
按题给定义积分：
\[
\tau_0=\int_0^\tau
\sqrt{1-v(t)^2/c^2}\,\mathrm dt
=\int_0^\tau\frac{\mathrm dt}
{\sqrt{1+(at/c)^2}}.
\]
令 $u=at/c$，
利用 $\int\mathrm du/\sqrt{1+u^2}
=\operatorname{arsinh}u$，得
\[
\boxed{\tau_0=\frac ca
\operatorname{arsinh}\frac{a\tau}{c}
=\frac ca\ln\left[
\frac{a\tau}{c}+
\sqrt{1+\left(\frac{a\tau}{c}\right)^2}\right]}.
\]
代入 $a=10g$ 及地球年长，
$\boxed{\tau_0\simeq3.5\ \text{个月}}$。
\excheck 固有时间小于地球坐标时间；
积分不能把火箭的运动看作恒速时间膨胀。

\exproblem{368}{59}
一粒子速率比光速小 $0.010\%$，
求其所谓“相对论质量”
$m=\gamma m_0$ 与静质量 $m_0$ 之比。
\exsolution
速度亏缺分数
$\delta=(c-v)/c=0.010\%=10^{-4}$，
故 $\beta=v/c=1-\delta=0.9999$。
根据定义
\[
\frac{m}{m_0}=\gamma
=\frac1{\sqrt{1-\beta^2}}
=\frac1{\sqrt{1-(1-\delta)^2}}
=\frac1{\sqrt{\delta(2-\delta)}}.
\]
代入得
\[
\boxed{\frac{m}{m_0}\simeq70.7\approx70}.
\]
\excheck 现代物理通常优先采用\emph{不变静质量}
$m_0$，把 $\gamma$ 归于动量和能量，
而不是把惯性质量随速度增长当成普适基本量。
此处沿用原书术语并显式解释。

\exproblem{369}{60}
一物体在自身静止系中的质量密度为
$\rho_0$。
求在哪个相对该物体运动的惯性系中，
测得密度比 $\rho_0$ 大
$\eta=25\%$。
\exsolution
题目所指质量密度沿用原书
\emph{静质量除以同时测得的空间体积}的定义。
运动方向的线性尺寸收缩为
$L=L_0/\gamma$，横截面积不变，
故同一组静质量所占空间体积
$V=V_0/\gamma$。
于是
\[
\rho=\frac{m_0}{V}
=\gamma\frac{m_0}{V_0}=\gamma\rho_0.
\]
令 $\rho/\rho_0=1+\eta$，
$\gamma=1+\eta$，
\[
\boxed{v=c\sqrt{1-\frac1{(1+\eta)^2}}
=c\frac{\sqrt{\eta(2+\eta)}}{1+\eta}
=0.60c}.
\]
\excheck 这是\emph{静质量/收缩体积}的密度，
并不是包含运动动能的能量密度除以 $c^2$；
不同密度定义不可混用。

\exproblem{370}{60}
一颗质子的相对论动量为
$p=10.0\,\mathrm{GeV}/c$。
求其速率比光速小的百分比。
\exsolution
质子静止能量
$m_pc^2\simeq0.938\,\mathrm{GeV}$。
由质能关系
\[
E^2=p^2c^2+m_p^2c^4,\qquad v=\frac{pc^2}{E},
\]
得
\[
\beta=\frac vc
=\frac{pc}{\sqrt{(pc)^2+(m_pc^2)^2}}
=\left[1+\left(\frac{m_pc}{p}\right)^2\right]^{-1/2}.
\]
因此相对光速亏缺的百分比为
\[
\boxed{100(1-\beta)=100
\left[1-
\frac{1}{\sqrt{1+(0.938/10.0)^2}}\right]
\simeq0.44\%}.
\]
\excheck 高速下质子动量不能近似写作 $m_pv$；
此处用的是完整相对论能量—动量关系。
